\scp{Wallacean subgroups in eastern Indonesia and Timor-Leste}{01-09}
In spite of all the issues muddying the waters described in previous sections
and in the next chapter,
there are nevertheless some patterns to be found.
The results are summarised here,
and then the patterns are detailed
and discussed in the relevant chapters.

The Wallacean subgroups presented below should be taken as preliminary,
partly because good descriptive work is so sparse,
and also because many previous claims are not supported on closer inspection,
so we are forging new ground.
We do not present this classification as definitive
and we fully expect that further work based on more complete
and more reliable data will lead to many refinements, improvements and corrections.
We have identified a number of areas where further work is needed.

We provide a top-down approach to subgrouping in this region looking at the historical phonologies.
We compare PMP etyma with their reflexes in identifiable cognates
and identify the sound changes which have occurred between the two.
While this approach goes far beyond the lexical classification which has often been the main approach
to language relatedness in this region,
it has its limitations.
Most importantly it ignores many facts revealed by a bottom-up analysis of the data.
Such facts include, for example,
the reconstruction of prenasalised plosives in subgroups
and the reintroduction of /ŋ/ after previous loss.
A bottom-up approach to the data can also reveal patterns of relatedness
which are masked by a purely
top-down approach (see \citealp{Edwards2018a} and \citealp{Edwards2018c}).

All data are not equal.
A holistic approach to subgrouping would triangulate phonological,
morphosyntactic, and lexical evidence.
We have taken lexical similarity into account in our subgrouping,
though we privilege phonological evidence where the two disagree.
Similarly, where data are available,
morphology has been taken into account.
However, for many languages no morphological data is available.
Innovations that are not exclusive,
not replacements, and are found only sporadically are worth noting,
but do not constitute subgrouping arguments unless they converge with stronger arguments.
Some retentions are noteworthy,
not as a basis for subgrouping,
but to show contrast with other subgroups that do something different.

To date, we have not found anything in the historical phonology
to be able to link these subgroups together
at a level below PMP.
In chapters \ref{ch:OthSub} and \ref{ch:Mar}, we review various suggestions
and possibilities for linking languages in this region together at a higher level,
and we revisit some of the phonological
and lexical evidence for CEMP,
and discuss their implications in \crf{ch:MorLes}.

Discussion about each subgroup,
claims about them, and summary classification follows in the relevant chapters.
Splits, mergers and regular sound changes of PMP proto-phonemes are indicated wherever known.
See the relevant chapters for classification outlines expanded beyond the three levels below.

In the following outline, as throughout the book,
\tsc{small caps} indicates a node with languages or other nodes below it.
Regular font is used for language names.
See the List of abbreviations
and conventions beginning on page \pageref{ch:Abb}.
See \frf{01-06} for the location of the Wallacean subgroups summarised below.
Which of these subgroups
and which of their sub-nodes should be considered linkages,
and which should be considered subgroups in the classical sense is a task for future work.

%\begin{xltabular}{\textwidth}[l]{l>{\raggedright\arraybackslash}p{5cm}>{\raggedright\arraybackslash}X}
\begin{xltabular}{\textwidth}[l]{lllX}
\mc{4}{l}{\tsc{\tbf{Malayo-Polynesian}} (Wallacean subgroups)}							\\
\rcl\tbf{1.}&\mc{3}{l}{\tsc{\tbf{\ili{Bima}-Lembata}} (± 58--83 languages/varieties)}\\
&&\mc{2}{l}{\tsc{East Sumbawa}}\\
&&&\ili{Sanggar}\\
&&&\ili{Kolo}\\
&&&\ili{Bima}\\
&&\mc{2}{l}{\tsc{Flores}}\\
&&&\tsc{West \ili{Central} Flores}\\
&&&\tsc{Flores-Lembata}\\
&&\mc{2}{l}{\tsc{Sumba-Havu}}\\
&&&\tsc{Sumba}\\
&&&\tsc{\ili{Havu}-Dhao}\\
\rcl\tbf{2.}&\mc{3}{l}{\tsc{\tbf{Timor-Babar}} (± 64--75 languages/varieties)}\\
&&\mc{2}{l}{\tsc{Helong}}\\
&&&\ili{Funai} \ili{Helong}\\
&&&\tsc{Nuclear Helong}\\
&&\mc{2}{l}{\tsc{Rote-Meto}}\\
&&&\tsc{Nuclear Rote}\\
&&&\tsc{West Rote-Meto}\\
&&\mc{2}{l}{\tsc{Eastern Timor}}\\
&&&\ili{Tetun}\\
&&&\ili{Habun}\\
&&&\tsc{\ili{Midiki}-Waima{\Q}a}\\
&&\mc{2}{l}{\tsc{\ili{Lakalei}-Idate}}\\
&&&\ili{Lakalei}\\
&&&\ili{Idate}\\
&\\
&&\mc{2}{l}{\tsc{Wetar-Atauro}}\\
&&&\ili{Galolen}\\
&&&\tsc{Wetar}\\
&&&\tsc{Atauro}\\
&&\mc{2}{l}{\tsc{Southwest Maluku}}\\
&&&\ili{Makuva}\\
&&&\tsc{\ili{Kisar}-Luang}\\
&&&East \ili{Damer}\\
&&&\tsc{\ili{Teun}-\ili{Nila}-Serua}\\
&&&\tsc{Babar}\\
&&\mc{2}{l}{\tsc{Southwest Tanimbar}}\\
&&&\ili{Selaru}\\
&&&\tsc{\ili{Seluwasan}-Makatian}\\
\rcl\tbf{3.}&\mc{3}{l}{\tsc{\tbf{\ili{Central} Timor}} (± 6 languages)}\\
&&\mc{2}{l}{Welaun}\\
&&\mc{2}{l}{\tsc{Nuclear \ili{Central} Timor}}\\
&&&\ili{Kemak}\\
&&&\ili{Tokodede}\\
&&&\tsc{Mambae}\\
\rcl\tbf{4.}&\mc{3}{l}{\tsc{\tbf{Aru}} (± 17 languages)}\\
&&\mc{2}{l}{Ujir}\\
&&\mc{2}{l}{\tsc{Nuclear Aru}}\\
&&&\ili{Kompane}\\
&&&\ili{Kola}\\
&&&\ili{Lola}\\
&&&\tsc{South Aru}\\
\rcl\tbf{5.}&\mc{3}{l}{\tsc{\tbf{Taliabu}} (\tsc{\tbf{Eastern Celebic, Celebic}}) (3 languages)}\\
&&\mc{2}{l}{Taliabu}\\
&&\mc{2}{l}{Soboyo}\\
&&\mc{2}{l}{Kadai}\\
\rcl\tbf{6.}&\mc{3}{l}{\tsc{\tbf{Sula-Buru}} (5 languages)}\\
&&\mc{2}{l}{Ambelau}\\
&&\mc{2}{l}{\tsc{Nuclear Sula-Buru}}\\
&&&\tsc{Sula}\\
&&&\tsc{Buru}\\
&\\
&\\
&\\
&\\
\rcl\tbf{7.}&\mc{3}{l}{\tsc{\tbf{\ili{Ambon}-Seram}} (±62 languages/varieties)}\\
&&\mc{2}{l}{Boano}\\
&&\mc{2}{l}{\tsc{Nuclear \ili{Ambon}-Seram}}\\
&&&\tsc{\ili{Ambon}-West Seram}\\
&&&\tsc{Atamanu}\\
&&&\tsc{Patakai-Manusela}\\
&&\mc{2}{l}{\tsc{\ili{Seti}-Benggoi}}\\
&&&\tsc{Seti}\\
&&&\ili{Benggoi}\\
&&\mc{2}{l}{Hoti}\\
&&\mc{2}{l}{\ili{Salas} (classification uncertain)}\\
\rcl\tbf{8.}&\mc{3}{l}{\tsc{\tbf{Seram-Tanimbar-Bomberai}} (±22 languages/varieties)}\\
&&\mc{2}{l}{\tsc{East Seram}}\\
&&\mc{2}{l}{\tsc{Eastern Islands}}\\
&&&\tsc{Nuclear Eastern Islands}\\
&&&\ili{Kowiai}\\
&&\mc{2}{l}{\tsc{\ili{Teor}-Kur}}\\
&&&\ili{Teor}\\
&&&\ili{Kur}\\
&&\mc{2}{l}{\tsc{Tanimbar-Bomberai}}\\
&&&\ili{Yamdena}\\
&&&\tsc{Nuclear Tanimbar-Bomberai}\\
&&&\tsc{East Bomberai}\\
\rcl\tbf{9.}&\mc{3}{l}{\tbf{Banda} Isolate}\\
&&\mc{2}{l}{\ili{Banda}-Elat}\\
&&\mc{2}{l}{\ili{Banda}-Eli}\\
\rcl\tbf{10.}&\mc{3}{l}{\tsc{\tbf{Eastern Malayo-Polynesian}} \citep{Blust1978,Blust2013} }\\
\rcl&\mc{3}{l}{(EMP is beyond the scope of this book)}\\
&&\mc{2}{l}{\tsc{South Halmahera-West New Guinea} \citep{Kamholz2014}}\\
&&\mc{2}{l}{\tsc{Oceanic}}\\
\end{xltabular}

There are also contact creoles in the region that
do not fit normal historical-comparative genealogical
classification or normal patterns of intergenerational transmission of language.
They each have native language speakers (L1),
as well as second language speakers (L2).
They should not be considered “dialects” of standard languages.

\newpage
\begin{xltabular}{\textwidth}[l]{lX}
\rcl \mc{2}{l}{\tbf{Contact Creole, \ili{Malay}-based}} \\
&	\ili{Larantuka} \ili{Malay} [lrt] \\
&	\ili{Alor} \ili{Malay} \\
&	\ili{Kupang} \ili{Malay} [mkn] \\
&	\ili{Ambon} \ili{Malay} [abs] \\
&	\ili{Banda} \ili{Malay} [bpq] \\
&	\ili{Aru} \ili{Malay} (a.k.a. Dobo \ili{Malay}) \\
\rcl \mc{2}{l}{\tbf{Contact Creole, \ili{Tetun}-based}} \\
&	\ili{Tetun} Dili [tdt] \\
\end{xltabular}


%\yytree{
%\yy[1,5,8]{\tsc{\tbf{1. \ili{Bima}-Lembata}} (± 58--83 languages/varieties)}			\\
%&	\yy[1,2,3]{\tsc{East Sumbawa}}		\\
%&	&	\yy{Sanggar}	\\
%&	&	\yy{Bima}	\\
%&	&	\yy{\tsc{\ili{Kolo}-Toloweri}}	\\
%&	\yy[1,2]{\tsc{Flores}}		\\
%&	&	\yy{\tsc{West \ili{Central} Flores}}	\\
%&	&	\yy{\tsc{Flores-Lembata}}	\\
%&	\yy[1,2]{\tsc{Sumba-Havu}}		\\
%&	&	\yy{\tsc{Sumba}}	\\
%&	&	\yy{\tsc{\ili{Havu}-Dhao}}	\\
%}
%
%\yytree{
%\yy[1,4,7,11,14,18,24]{\tsc{\tbf{2. Timor-Babar}} (± 64--75 languages/varieties)}			\\
%&	\yy[1,2]{\tsc{Helong}}		\\
%&	&	\yy{\ili{Funai} Helong}	\\
%&	&	\yy{\tsc{Nuclear Helong}}	\\
%&	\yy[1,2]{\tsc{Rote-Meto}}		\\
%&	&	\yy{\tsc{Nuclear Rote}}	\\
%&	&	\yy{\tsc{West Rote-Meto}}	\\
%&	\yy[1,2,3]{\tsc{Eastern Timor}}		\\
%&	&	\yy{Tetun}	\\
%&	&	\yy{Habun}	\\
%&	&	\yy{\tsc{\ili{Midiki}-Waima{\Q}a}}	\\
%&	\yy[1,2]{\tsc{\ili{Lakalei}-Idate}}		\\
%&	&	\yy{Lakalei}	\\
%&	&	\yy{Idate}	\\
%&	\yy[1,2,3]{\tsc{Wetar-Atauro}}		\\
%&	&	\yy{Galolen}	\\
%&	&	\yy{\tsc{Wetar}}	\\
%&	&	\yy{\tsc{Atauro}}	\\
%&	\yy[1,2,3,4,5]{\tsc{Southwest Maluku}}		\\
%&	&	\yy{Makuva}	\\
%&	&	\yy{\tsc{\ili{Kisar}-Luang}}	\\
%&	&	\yy{East Damer}	\\
%&	&	\yy{\tsc{\ili{Teun}-\ili{Nila}-Serua}}	\\
%&	&	\yy{\tsc{Babar}}	\\
%&	\yy[1,2]{\tsc{Southwest Tanimbar}}		\\
%&	&	\yy{Selaru}	\\
%&	&	\yy{\tsc{\ili{Seluwasan}-Makatian}}	\\
%}
